\scanpage{295}
\mspage{211}

Le ``il faut'' \unsure{est} \unsure{remis|mis} pour mémoire, on vient de le voir comme conséquence du fait que tout objet de $M$ est accessible. Le ``il suffit'' est essentiellement SGA 4 I 8.13.2. b), sauf que dans loc. cit on suppose $M$ \unsure{muni|munie} d'une \unsure{fc} \uline{stricte}. Mais la démonstration donnée dans loc. cit. marche aussi pour une \unsure{fcga}, en choisissant dans loc cit (p.137) $\pi$ tel que $\unsure{\xi|\zeta}$ soit $\pi$-accessible, et de plus $\pi$ soit \uline{\unsure{utile|stable}} pour la filtration cardinale gén. envisagée.

Reste à prouver $c''\overset{?}{\Longrightarrow}c$, sous réserve que ce soit bien vrai. Ce serait une conséquence du lemme suivant (où $D,\pi$ sont \unsure{remplacés|remplis} par $C,\pi_0$)

\uline{Lemme} 4.8.12 \unsure{(?)}. Soient $C$ une petite catégorie,

\margin{Lemme non prouvé \unsure{ici}, cf dém 4.11. (\unsure{en addenda|en admissib.})}% slanted

$\pi_0$ un cardinal infini. Si $C$ est stable par $\varprojlim$ finies, les $\operatorname{Ind}_{\pi_0}C$ également.

L'assertion similaire pour $\operatorname{Ind}(C)$, \unsure{ald|a.ld} $\operatorname{Ind}_{\pi}C$, est établie dans SGA 4 I 8.9.1 c). Mais il n'est pas clair sans plus que $\operatorname{Ind}_{\pi}C$ soit stable dans $\operatorname{Ind}(C)$ par $\varprojlim$ finies, dans le cas envisagé. \deleted{C'est lié à la} \unsure{D'en|Pour} \unsure{faire|bien} préciser, si $F=\varprojlim_{j\in J}F_j$, les $F_j$ dans $\operatorname{Ind}_{\pi}C$,
% archive sheet number: 294
